\chapter{Experimentación}

\section{Escenario Base (Baseline)}

El escenario base corresponde a la ejecución del sistema con configuración funcional mínima por control (sin endurecimiento adicional), permitiendo establecer una línea de referencia para la comparación de resultados en cada uno de los controles evaluados (C1--C4).

En la matriz experimental, este escenario se representa mediante las configuraciones C1-B, C2-B, C3-B y C4-B, evaluadas en los niveles de carga VUS \{1, 5, 10, 20\}.

\begin{figure}[H]
    \centering
    \caption{Propuesta de arquitectura del escenario base}
    \includegraphics[width=0.7\textwidth]{Pictures/arquitectura/DiagramaBase.drawio.png}
    \label{fig:DiagramaBase}

    \vspace{0.2cm}
    {\raggedright \small \textit{Nota.} Arquitectura base de referencia para comparación entre variantes de control. \par}
\end{figure}
\FloatBarrier
\subsection{Control 1: API Gateway}

Este control responde al problema de gobernar el tráfico de entrada (\textit{north-south}) hacia la arquitectura de microservicios.

\begin{figure}[H]
    \centering
    \caption{Propuesta arquitectura escenario C1}
    \includegraphics[width=0.7\textwidth]{Pictures/arquitectura/C1.drawio.png}
    \label{fig:C1}

    \vspace{0.2cm}
    {\raggedright \small \textit{Nota.} C1. \par}
\end{figure}

En la matriz experimental, se definen los siguientes escenarios:

\begin{itemize}
    \item C1-B: baseline con NGINX Ingress como ruta de entrada estándar.
    \item C1-A: implementación con Istio Gateway.
    \item C1-C: implementación con Kong API Gateway.
\end{itemize}

\textbf{Trade-off:} simplicidad y rendimiento versus capacidades avanzadas de gobierno y seguridad.

En este contexto, las variantes comparadas son equivalentes en alcance funcional porque operan sobre el mismo punto de control de borde (\textit{edge}), donde se gestiona el tráfico externo antes de su ingreso al dominio interno de microservicios.

%------------------------------------------------

\section{Control 2: mTLS mediante Service Mesh}

Este control permite asegurar la comunicación entre microservicios mediante cifrado en tránsito y autenticación mutua.

\begin{figure}[H]
    \centering
    \caption{Propuesta arquitectura escenario C2}
    \includegraphics[width=0.7\textwidth]{Pictures/arquitectura/C2.drawio.png}
    \label{fig:C2}

    \vspace{0.2cm}
    {\raggedright \small \textit{Nota.} C2. \par}
\end{figure}

Se consideran tecnologías como Istio y Linkerd, las cuales implementan mTLS de forma transparente para los servicios de aplicación.

Este control impacta directamente la confidencialidad, integridad y autenticidad de las comunicaciones internas, aunque introduce sobrecarga en latencia y consumo de recursos.

En la matriz experimental, se definen los siguientes escenarios:

\begin{itemize}
    \item C2-B: escenario base sin mTLS.
    \item C2-A: implementación con Istio mTLS.
    \item C2-C: implementación con Linkerd mTLS.
\end{itemize}

\textbf{Trade-off:} mayor garantía criptográfica entre servicios versus incremento de complejidad operativa y consumo de recursos.

%------------------------------------------------

\section{Control 3: Network Policies}

Las políticas de red permiten implementar microsegmentación dentro del clúster de Kubernetes, restringiendo la comunicación entre servicios únicamente a los flujos necesarios.

\begin{figure}[H]
    \centering
    \caption{Propuesta arquitectura escenario C3}
    \includegraphics[width=0.7\textwidth]{Pictures/arquitectura/C3.drawio.png}
    \label{fig:C3}

    \vspace{0.2cm}
    {\raggedright \small \textit{Nota.} C3. \par}
\end{figure}

Este control reduce la superficie de ataque y limita la propagación lateral de amenazas dentro del sistema.

Su implementación se realiza mediante \texttt{NetworkPolicy} de Kubernetes en perfiles \textit{basic} y \textit{strict}.

Este control aborda la gestión del tráfico este-oeste dentro del clúster.

En la matriz experimental, se definen los siguientes escenarios:

\begin{itemize}
    \item C3-B: escenario base sin políticas de red.
    \item C3-A: segmentación básica.
    \item C3-C: segmentación estricta.
\end{itemize}

\textbf{Trade-off:} mayor aislamiento de red y reducción de superficie de ataque versus riesgo de afectar conectividad funcional en perfiles más restrictivos.


%------------------------------------------------

\section{Control 4: Rate Limiting}

El control de limitación de solicitudes permite regular el volumen de tráfico aceptado por los servicios, mitigando ataques como denegación de servicio (DoS) y abuso de APIs.

\begin{figure}[H]
    \centering
    \caption{Propuesta arquitectura escenario C4}
    \includegraphics[width=0.7\textwidth]{Pictures/arquitectura/C4.drawio.png}
    \label{fig:C4}

    \vspace{0.2cm}
    {\raggedright \small \textit{Nota.} C4. \par}
\end{figure}

En la matriz experimental, se definen los siguientes escenarios:

\begin{itemize}
    \item C4-B: escenario base sin limitación.
    \item C4-A: limitación moderada (1200 rpm).
    \item C4-C: limitación estricta (300 rpm).
\end{itemize}

Este control destaca por su bajo costo de implementación y su impacto en la estabilidad del sistema bajo cargas elevadas, aunque su contribución a la seguridad estructural es más limitada en comparación con otros controles. En la ejecución experimental, la limitación se aplicó en ingress mediante anotaciones de NGINX.

Para mantener la comparabilidad metodológica, las variantes se evaluaron dentro de cada control (C1, C2, C3 y C4) bajo los mismos niveles de carga definidos en la matriz experimental. En particular, C1 se evaluó sobre la ruta de entrada de gateway/ingress, mientras que C2--C4 se ejecutaron sobre la ruta de servicios internos expuesta para benchmark; por ello, las conclusiones se interpretan de forma intra-control y no como comparación directa absoluta entre controles diferentes.

%------------------------------------------------
